\documentclass[11pt]{article}
\usepackage{setspace}
\usepackage{geometry}
\geometry{margin=1in}
\setstretch{1.15}
\usepackage[hidelinks]{hyperref}

\begin{document}

% --- Applicant Information ---
\noindent
Zhongli Wang \\
Bielefeld, Germany \\
Phone: +49 160 9541 6015 \\
Email: nkwangzhongli@gmail.com \\
Website: \href{https://sites.google.com/view/zhongli-wang}{https://sites.google.com/view/zhongli-wang}

\vspace{0.45cm}

% --- Recipient Information ---
\noindent
Hiring Committee \\
Transforming Economies Lab (IT:U) \& Transforming Economies Group (CSH) \\
Linz, Austria

\vspace{0.45cm}

% --- Subject Line ---
\noindent
Application for Postdoctoral Researcher --- Digital Traces of Software Development (IT:U--CSH)

\vspace{0.35cm}

\noindent Dear Members of the Hiring Committee,\\[0.2cm]

I am writing to apply for the Postdoctoral Researcher position in the Transforming Economies Lab at IT:U, in collaboration with the Transforming Economies Group at the Complexity Science Hub. I am completing my Ph.D.\ in Economics at Bielefeld University and Universit\`a Cattolica del Sacro Cuore through the Marie Sk{\l}odowska--Curie EPOC programme. My research combines applied econometrics, computational social science, and scalable simulation and data workflows to study how technology shapes behavior, productivity, and market outcomes. 
Your project---building high-resolution data pipelines from digital traces of software development and using them to study labor markets, geography, careers, and the role of automation and AI---is a close match to both my skill set and my research ambitions.

My scientific vision is to develop a measurement-driven research program on \emph{digital production and labor-market organization} in the era of AI-assisted work. 
I am particularly interested in how software developers' careers and collaboration patterns respond to changing technologies, and how these responses differ across places, firms, and teams. Substantively, I aim to connect micro-level work traces to outcomes such as mobility, specialization, human-capital accumulation, team performance, and local labor-market structure. Methodologically, I want to combine (i) trace-based measures of effort, learning, and collaboration from line-by-line activity and repository metadata with (ii) econometric designs that support credible inference about technology shocks and organizational change. The joint IT:U--CSH environment is especially attractive to me because it couples infrastructure development with substantive inquiry, enabling research that is both reproducible and genuinely data-rich.

A core strength I would bring to the project is experience designing and maintaining scalable computational workflows. In my dissertation work, I developed large-scale simulation pipelines (including parallel execution on high-performance computing clusters) to analyze strategic behavior and heterogeneous decision rules in data-rich markets. This required building modular codebases, managing large outputs, and implementing reproducible experiment tracking, parameter sweeps, and quality checks. I am comfortable working with Python-based workflows, performance constraints, and version-controlled research pipelines. While the data modality in your project differs from simulation output, the engineering principles are the same: transforming raw, high-frequency records into validated, well-documented datasets suitable for research and collaboration. I would be excited to contribute to the design and extension of your pipeline architecture, including data validation, feature extraction, and workflow automation.
% that supports both the team and future users.

I also bring strong training in quantitative empirical methods and a clear plan for how I would use them in this setting. My applied toolkit includes regression-based analysis for panel and repeated-observation settings, event-study style designs for technology adoption and policy or platform changes, and careful robustness strategies for measurement uncertainty. 
% Where identification is feasible, I am interested in implementing quasi-experimental approaches (e.g., staggered adoption designs, difference-in-differences with appropriate diagnostics, or instrumented exposure measures) to separate the effects of new tools from confounders such as selection and endogenous team assignment. 
I am also comfortable combining econometrics with predictive components when useful (e.g., using machine learning to construct high-dimensional measures from traces, while keeping the causal estimand and inference transparent). 
% This matches your emphasis on rigorous quantitative analysis, including causal identification ``where appropriate,'' while remaining grounded in the realities of digital trace data.
The project topics---local labor markets in software development, spatial structure, and career trajectories---connect naturally to my broader interests in economic dynamics and technology-mediated behavior. Conceptually, the software labor market is a leading case for understanding modern skill formation and task reallocation: work is partially observable, collaboration is digitally logged, and the adoption of automation can be measured at scale. I am especially interested in questions such as: How does AI-tool adoption change the within-person evolution of productivity and specialization? Does automation alter the geography of opportunity by shifting which tasks can be done remotely or by reducing the returns to agglomeration? 
% Which teams benefit most, and do these changes widen or narrow inequality across workers and regions? 
These questions align closely with the Transforming Economies program and are also highly policy-relevant, with implications for training, innovation policy, and regional development.

My research style is collaborative and interdisciplinary. During my doctoral training, I worked in international teams and coauthored projects that sit at the intersection of economics and modern computational methods. One project develops an \emph{agentic workflow} that integrates large-language-model agents into research pipelines, with an emphasis on reproducibility, documentation, and modular task delegation (e.g., data collection, coding support, and structured reporting). This experience has shaped how I build research infrastructure: I prioritize transparent logging, human-auditable outputs, and workflows that lower the cost of collaboration. 
% I would welcome the opportunity to contribute to shared infrastructure at IT:U--CSH while also leading and co-developing empirical research with colleagues from economics, computational social science, and complexity research.

Regarding training goals, I see this position as an ideal opportunity to deepen two complementary competencies while maintaining a strong empirical profile. First, I want to further professionalize my data engineering skills around high-resolution trace data: efficient storage, incremental processing, robust feature extraction, and pipeline testing. Second, I want to expand my applied econometric practice specifically for digital traces: measurement models for work intensity and learning, linking trace-derived metrics to external outcomes, and careful identification strategies around technology adoption. 
% The IT:U--CSH joint environment, with access to high-performance computing resources and a culture of methodological innovation, is exactly the setting where these goals are realistic.

% Finally, I can confirm that I am willing and able to be present onsite in Vienna and/or Linz as required. I value being embedded in the research environment through seminars, workshops, and day-to-day collaboration, and I would be enthusiastic about contributing to the interdisciplinary research life of both IT:U and CSH. I am also happy to participate in light project-based teaching and supervision activities and to help build a supportive environment for students and junior researchers.

Thank you for considering my application! I would welcome the opportunity to discuss how my background in quantitative empirical analysis, scalable computational workflows, and research on technology-driven economic behavior can contribute to the Transforming Economies Lab/Group.\\[0.35cm]

\noindent Sincerely,\\[0.9cm]
\noindent Zhongli Wang

\end{document}